\section*{Capítulo \ref{ch:g4:adapted-to-their-world} --- Adaptados a su mundo}

\begin{solution}{exo:g4:adapted-to-their-world:1}
Un rasgo del cuerpo o del comportamiento que ajusta a un \omterm{def:g1:living-or-not:living}{ser vivo} a
las condiciones de su \omterm{def:g2:caring-for-nature:environment}{entorno}. \omterm{def:g1:animals-around-us:animal}{Animal}: el doble aislamiento de \omterm{def:g1:animals-around-us:covering}{pelo} y
grasa del oso polar. \omterm{def:g1:plants-around-us:plant}{Planta}: el \omterm{def:g1:plants-around-us:parts}{tallo} del cactus que almacena agua.
\end{solution}

\begin{solution}{exo:g4:adapted-to-their-world:2}
Nada de agua: beber cien litros de golpe y aguantar días sin beber,
más la despensa de grasa de las jorobas. Arena blanda: pies anchos de
dos dedos. Polvo levantado por el viento: ollares que se cierran y
pestañas dobles y largas. (Calor y hambre: otra vez las jorobas ---
grasa para las travesías flacas.)
\end{solution}

\begin{solution}{exo:g4:adapted-to-their-world:3}
Frío: \omterm{def:g1:animals-around-us:covering}{pelo} espeso sobre una capa profunda de grasa. Agua helada y
nieve: zarpas anchas y algo palmeadas --- raquetas de nieve y remos.
Presas desconfiadas sobre hielo blanco: el camuflaje de un pelaje
blanco.
\end{solution}

\begin{solution}{exo:g4:adapted-to-their-world:4}
Las púas son \omterm{def:g1:plants-around-us:parts}{hojas} reducidas: casi no pierden agua y defienden a la
\omterm{def:g1:plants-around-us:plant}{planta} de los \omterm{def:g1:animals-around-us:animal}{animales} sedientos. El \omterm{def:g1:plants-around-us:parts}{tallo} en forma de barril almacena
agua y ha asumido el trabajo de atrapar la luz que harían las \omterm{def:g1:plants-around-us:parts}{hojas}.
\end{solution}

\begin{solution}{exo:g4:adapted-to-their-world:5}
La forma de cojín rompe el viento helado: dentro de él el aire se
queda más templado. La forma misma es el refugio.
\end{solution}

\begin{solution}{exo:g4:adapted-to-their-world:6}
Cazados: colores de camuflaje, armadura o púas, velocidad, \omterm{def:g1:the-five-senses:sense}{ojos} muy
separados que vigilan alrededor (con dos basta). Cazadores: \omterm{def:g1:the-five-senses:sense}{ojos} hacia
delante, sentidos finos, garras, velocidad (con dos basta).
\end{solution}

\begin{solution}{exo:g4:adapted-to-their-world:7}
\omterm{def:g1:animals-around-us:covering}{Pelo} denso y suave --- un mamífero; enormes garras delanteras en forma
de pala --- un excavador; \omterm{def:g1:the-five-senses:sense}{ojos} diminutos y débiles --- vida en la
oscuridad. Todo junto: un mamífero excavador subterráneo, el topo.
\end{solution}

\begin{solution}{exo:g4:adapted-to-their-world:8}
Respuesta abierta. Ejemplo del pato: pies palmeados (remar), \omterm{def:g1:animals-around-us:covering}{plumas}
que repelen el agua (mantenerse seco y caliente a flote), pico plano
de cribar (comer del agua).
\end{solution}

\begin{solution}{exo:g4:adapted-to-their-world:9}
Las \omterm{def:g1:the-five-senses:sense}{orejas} son radiadores: las grandes y finas sueltan el calor del
cuerpo. En el desierto, soltar calor es el problema que hay que
resolver --- unas \omterm{def:g1:the-five-senses:sense}{orejas} enormes refrescan al zorro. En el hielo,
conservar el calor lo es todo --- unas \omterm{def:g1:the-five-senses:sense}{orejas} diminutas pierden lo
mínimo. Un mismo órgano, dos climas, dos tamaños contrarios.
\end{solution}

\begin{solution}{exo:g4:adapted-to-their-world:10}
El cuerpo afinado, los pies palmeados y aguantar la \omterm{def:g4:breathing:movements}{respiración} dicen:
un nadador y buceador; el \omterm{def:g1:animals-around-us:covering}{pelo} denso que repele el agua dice: un
mamífero que vive medio metido en agua fría; los bigotes dicen: tantear
el camino en aguas oscuras o turbias. Un mamífero acuático constructor
de presas, de ríos y lagos --- el castor.
\end{solution}

\begin{solution}{exo:g4:adapted-to-their-world:11}
Lo impuso el agua. La ballena y la trucha no heredaron nada con forma
de nadador de una familia común --- una es un mamífero y la otra un
\omterm{prop:g3:sorting-living-things:groups}{pez} --- y sin embargo el mismo problema, moverse deprisa por el agua,
favorece en las dos la misma respuesta afinada. Las adaptaciones
responden a los problemas del \omterm{def:g2:caring-for-nature:environment}{entorno}; problemas parecidos pueden
sacar formas parecidas de familias muy distintas.
\end{solution}
